% ============================================================================
% CUSTOM COMMANDS & NOTATIONS
% ============================================================================
% Define commonly used mathematical notation and utility commands
% Include this file in wiki_main.tex or individual sections
% ============================================================================

% ============================================================================
% COSMOLOGY & ASTROPHYSICS
% ============================================================================

% Cosmological parameters
\newcommand{\Om}{\Omega_{\mathrm{m}}}           % Matter density parameter
\newcommand{\OL}{\Omega_{\Lambda}}              % Dark energy density parameter
\newcommand{\Ob}{\Omega_{\mathrm{b}}}           % Baryon density parameter
\newcommand{\Ok}{\Omega_{\mathrm{k}}}           % Spatial curvature parameter
\renewcommand{\H}[1][]{H_{0#1}}                % Hubble constant (optional redshift arg)
\newcommand{\Hz}[1]{H\left(#1\right)}          % H(z) function

% Distances & coordinates
% Keep \chi as standard LaTeX Greek letter
\newcommand{\chiz}[1]{\chi\left(#1\right)}     % χ(z)
\newcommand{\ra}{\mathrm{r}}                   % Radial coordinate in cosmology
\newcommand{\Da}{D_A}                          % Angular diameter distance
\newcommand{\Dz}[1]{D_A\left(#1\right)}        % D_A(z)

% Growth factor & perturbations
\newcommand{\Dg}{D_+}                          % Growth factor (linear)
\newcommand{\Dgz}[1]{D_+\left(#1\right)}       % D_+(z)
% \delta redefined as matter dens perturbation - keep standard LaTeX \delta
\newcommand{\deltam}{\delta_{\mathrm{m}}}      % Matter density perturbation
\newcommand{\deltamz}[2]{\delta_{\mathrm{m}}\left(#1, #2\right)}  % δ_m(k,z)

% ============================================================================
% POWER SPECTRUM & CORRELATION FUNCTIONS
% ============================================================================

% Cartesian Fourier
\newcommand{\Pk}{P(k)}                         % 3D power spectrum
\newcommand{\Pkz}[2]{P(k, #1, #2)}             % P(k, z₁, z₂)
\newcommand{\xik}{\xi(r)}                      % Correlation function
\newcommand{\xikz}[2]{\xi(r, #1, #2)}          % ξ(r, z₁, z₂)

% Spherical decomposition
\newcommand{\Clk}{C_{\ell}(k)}                 % SFB power spectrum
\newcommand{\Clkz}[2]{C_{\ell}(k, #1, #2)}     % C_ℓ(k, z₁, z₂)
\newcommand{\Cl}{C_{\ell}}                     % Angular power spectrum
\newcommand{\Clij}[2]{C^{AB}_{\ell #1, #2}}    % Cross-spectrum in bins

% ============================================================================
% WEAK LENSING
% ============================================================================

% Lensing quantities
% Keep standard LaTeX Greek letters \kappa, \phi, \Psi
\newcommand{\convergence}{\kappa}              % Convergence field
\newcommand{\ShearField}{\gamma}               % Shear magnitude (renamed)
\newcommand{\shearComplex}{{\gamma_1 + i\gamma_2}}   % Complex shear notation
\newcommand{\ShearE}{\gamma^E}                 % E-mode shear
\newcommand{\ShearB}{\gamma^B}                 % B-mode shear
\newcommand{\lensingpot}{\varphi}              % Lensing potential (typically φ or ψ)
\newcommand{\potpert}{\Psi}                    % Potential perturbation

% Lensing kernels
\newcommand{\Wkappa}{W^{\kappa}}               % Convergence kernel
\newcommand{\Wshear}{W^{\gamma}}               % Shear kernel
\newcommand{\Wden}{W^{\mathrm{den}}}           % Density kernel
\newcommand{\Wl}[1]{W^{#1}_{\ell}}             % Generic lensing kernel

% ============================================================================
% SPHERICAL FOURIER-BESSEL
% ============================================================================

% Bessel functions
\newcommand{\Jell}[1]{j_{\ell}\left(#1\right)}     % Spherical Bessel j_ℓ(x)
\newcommand{\Jnu}[1]{j_{\nu}\left(#1\right)}       % j_ν(x)
\newcommand{\Nell}[1]{n_{\ell}\left(#1\right)}     % Neumann function n_ℓ(x)
\newcommand{\jzero}[1]{j_0\left(#1\right)}        % j₀(x) - monopole
\newcommand{\jone}[1]{j_1\left(#1\right)}         % j₁(x) - dipole
\newcommand{\jtwo}[1]{j_2\left(#1\right)}         % j₂(x) - quadrupole

% Derivatives of Bessel functions
\newcommand{\Jdell}[1]{j'_{\ell}\left(#1\right)}   % dj_ℓ/dx
\newcommand{\Jdnu}[1]{j'_{\nu}\left(#1\right)}     % dj_ν/dx

% SFB coefficients & expansions
\newcommand{\flm}{f_{\ell m}}                  % SFB coefficient f_ℓm
\newcommand{\flmk}[1]{f_{\ell m}\left(#1\right)}  % f_ℓm(k)
\newcommand{\alphanm}{\alpha_{n\ell m}}        % α_{nℓm} basis coefficient
\newcommand{\alphank}[1]{\alpha_{n\ell}\left(#1\right)}  % α_{nℓ}(k)

% SFB grids
\newcommand{\rln}[1]{r^{(\ell)}_{#1}}          % r_ℓ,n - radial grid point
\newcommand{\xln}[1]{x^{(\ell)}_{#1}}          % x_ℓ,n - zero of Bessel function

% ============================================================================
% SPHERICAL HARMONICS
% ============================================================================

% Standard & spin-weighted
\newcommand{\YellM}{Y_{\ell m}}                % Standard spherical harmonic
\newcommand{\YellMarg}[1]{Y_{\ell m}\left(#1\right)}   % Y_ℓm(θ,φ)
\newcommand{\sYlm}[1]{{}_{#1}Y_{\ell m}}       % Spin-s harmonic ₛY_ℓm
\newcommand{\spmweighted}[1]{{}_{{\pm #1}}Y_{\ell m}}    % Spin-weight ± notation
\newcommand{\aelm}{a_{\ell m}}                 % Harmonic coefficient a_ℓm
\newcommand{\aelmAB}[2]{a^{#1#2}_{\ell m}}     % Cross-spectrum coeff a^AB_ℓm

% ============================================================================
% INTEGRATION & QUADRATURE
% ============================================================================

% Integral symbols with limits
\newcommand{\intk}{\int_0^\infty dk}           % Integral over wavenumber
\newcommand{\intz}{\int_0^\infty dz}           % Integral over redshift
\newcommand{\intchi}{\int_0^\infty d\chi}      % Integral over comoving distance
\newcommand{\intOmega}{\int d\Omega}           % Integral over solid angle

% Quadrature notation
\newcommand{\quadrule}[3]{\left\langle #1 \right\rangle_{#2}^{(#3)}}  % Quadrature rule

% ============================================================================
% OPERATORS
% ============================================================================

% Laplacian operators
\newcommand{\Lap}{\nabla^2}                    % Laplacian (3D)
\newcommand{\Lapperp}{\nabla^2_\perp}          % Laplacian on sphere
\newcommand{\Lapang}{\nabla^2_\Omega}          % Angular Laplacian

% Derivatives on sphere
\newcommand{\dtheta}{\partial_\theta}          % ∂/∂θ
\newcommand{\dphi}{\partial_\phi}              % ∂/∂φ

% ============================================================================
% PROBABILITY & STATISTICS
% ============================================================================

% Ensemble averages
\newcommand{\avg}[1]{\left\langle #1 \right\rangle}     % Ensemble average ⟨...⟩
\newcommand{\var}[1]{\text{Var}\left(#1\right)}         % Variance
\newcommand{\cov}[2]{\text{Cov}\left(#1, #2\right)}    % Covariance

% Delta functions
\newcommand{\deltaf}[1]{\delta\left(#1\right)}          % Dirac delta δ(x)
\newcommand{\deltathree}[1]{\delta^{(3)}\left(#1\right)}    % 3D delta δ³(k)
\newcommand{\KD}[2]{\delta_{#1 #2}}                     % Kronecker delta δ_ij

% ============================================================================
% PHYSICAL CONSTANTS & UNITS
% ============================================================================

\newcommand{\cc}{c}                            % Speed of light
\newcommand{\GG}{G}                            % Gravitational constant
\newcommand{\rhobar}{\bar{\rho}_m}             % Mean matter density
\newcommand{\rhocrit}{\rho_{\mathrm{crit}}}    % Critical density
\newcommand{\pcrit}{\rho_c}                    % Critical dens. (short form)

% ============================================================================
% NOTATION: SPECIAL VECTORS & MATRICES
% ============================================================================

% Vector notation (using bold)
\newcommand{\vk}{\mathbf{k}}                   % k-vector
\newcommand{\vr}{\mathbf{r}}                   % r-vector
\newcommand{\vx}{\mathbf{x}}                   % x-vector
\newcommand{\vtheta}{{\boldsymbol{\theta}}}    % Angular coordinates vector

% Matrix notation
\newcommand{\mat}[1]{\mathbf{#1}}              % Matrix notation

% ============================================================================
% TEXT FORMATTING SHORTCUTS
% ============================================================================

\newcommand{\SIH}{\textit{Statistical Isotropy and Homogeneity}}  % SIH
\newcommand{\limber}{\textit{Limber approximation}}               % Limber
\newcommand{\sfb}{\textit{Spherical Fourier-Bessel}}             % SFB

% ============================================================================
% BOXES & HIGHLIGHTS
% ============================================================================

% Important result box
\newcommand{\important}[1]{%
  \begin{tcolorbox}[colback=yellow!5,colframe=orange!75!black,title=Important]
    #1
  \end{tcolorbox}%
}

% Key equation box
\newcommand{\keyeq}[2]{%
  \begin{tcolorbox}[colback=blue!5,colframe=blue!75!black,title=Key Equation]
    \textbf{#1}: $#2$
  \end{tcolorbox}%
}

% Note box
\newcommand{\sidenote}[1]{%
  \begin{tcolorbox}[colback=gray!10,colframe=gray!75!black,title=Note]
    #1
  \end{tcolorbox}%
}

% ============================================================================
% END CUSTOM COMMANDS
% ============================================================================
